% estilo_diagramas.tex - Estilo comun de los cinco diagramas TikZ de la tesis
% (arquitectura, u1, u2, u3, u4). Cada diagrama es un documento standalone que hace
%   \documentclass[10pt,border=0pt]{standalone}
%   \usepackage[utf8]{inputenc}  \usepackage[T1]{fontenc}  \usepackage[spanish]{babel}
%   \usepackage{tikz}            % estilo_diagramas.tex - Estilo comun de los cinco diagramas TikZ de la tesis
% (arquitectura, u1, u2, u3, u4). Cada diagrama es un documento standalone que hace
%   \documentclass[10pt,border=0pt]{standalone}
%   \usepackage[utf8]{inputenc}  \usepackage[T1]{fontenc}  \usepackage[spanish]{babel}
%   \usepackage{tikz}            % estilo_diagramas.tex - Estilo comun de los cinco diagramas TikZ de la tesis
% (arquitectura, u1, u2, u3, u4). Cada diagrama es un documento standalone que hace
%   \documentclass[10pt,border=0pt]{standalone}
%   \usepackage[utf8]{inputenc}  \usepackage[T1]{fontenc}  \usepackage[spanish]{babel}
%   \usepackage{tikz}            % estilo_diagramas.tex - Estilo comun de los cinco diagramas TikZ de la tesis
% (arquitectura, u1, u2, u3, u4). Cada diagrama es un documento standalone que hace
%   \documentclass[10pt,border=0pt]{standalone}
%   \usepackage[utf8]{inputenc}  \usepackage[T1]{fontenc}  \usepackage[spanish]{babel}
%   \usepackage{tikz}            \input{estilo_diagramas}
% y dibuja a 15 cm exactos de ancho (\useasboundingbox), para que la tesis lo incluya con
% \includegraphics[width=\textwidth] a escala 1. Letra: \footnotesize (8 pt) en nodos y
% etiquetas, \small (9 pt) en los titulos de unidad. Nunca \scriptsize ni \tiny.
% construir.py compila <nombre>.tex y copia el PDF a ../images/diag_<nombre>.pdf.
% Un diagrama se compila aparte: dentro no van \ref ni \cite (se escribe "U2", "Anexo D").
%
% Lenguaje visual
%   color      = unidad a la que pertenece el elemento (U1 azul, U2 naranja, U3 verde,
%                U4 violeta; gris = fuera de las unidades); relleno muy claro y borde
%                oscuro, legible en blanco y negro. Un puerto lleva el color de la unidad
%                del otro lado de la interfaz (gris si son varias).
%   forma      = tipo de nodo: proceso (rectangulo redondeado), archivo (esquina cortada),
%                decision (rombo), inicio/fin (estadio), entorno SUMO + TraCI (gris, doble
%                borde), puerto de interfaz (flecha-etiqueta), semillas (pildora en cursiva).
%   trazo      = tipo de flecha: flujo (continua), interfaz I1-I8 (discontinua: dato de una
%                interfaz del catalogo, lleva \itf o \itfn), repeticion (continua gruesa),
%                no usado (punteada gris; el borde de un nodo previsto es igual).

\usetikzlibrary{positioning, fit, calc, arrows.meta, backgrounds, chains,
                shapes.geometric, shapes.misc, shapes.symbols, babel}

% ---------------------------------------------------------------- colores
% <u>b = borde y titulo; <u>f = fondo del grupo de la unidad; <u>s = fondo de un subgrupo
\definecolor{u1b}{HTML}{2F5597}  \definecolor{u1f}{HTML}{EDF2FA}  \definecolor{u1s}{HTML}{DCE6F5}
\definecolor{u2b}{HTML}{A65A0B}  \definecolor{u2f}{HTML}{FDF3E7}  \definecolor{u2s}{HTML}{F9E3CA}
\definecolor{u3b}{HTML}{3A6E21}  \definecolor{u3f}{HTML}{EEF6EA}  \definecolor{u3s}{HTML}{DDEDD5}
\definecolor{u4b}{HTML}{5E4A9A}  \definecolor{u4f}{HTML}{F2EFF9}  \definecolor{u4s}{HTML}{E3DDF2}
% externo: resultados de la tesis (Cap. V y Anexo A) y todo lo que no es de una unidad
\definecolor{extb}{HTML}{404040} \definecolor{extf}{HTML}{F2F2F2} \definecolor{exts}{HTML}{E6E6E6}
% unidad "actual" (ub, uf, us): la fija \unidad{u2} (o u1, u3, u4, ext) al comienzo de cada
% diagrama; los estilos sin argumento usan esa unidad y los que llevan =<u> usan la indicada
% (p. ej. proceso=u3 para un nodo de U3 dibujado en el diagrama de U2).
\colorlet{ub}{extb} \colorlet{uf}{extf} \colorlet{us}{exts}
\newcommand{\unidad}[1]{\colorlet{ub}{#1b}\colorlet{uf}{#1f}\colorlet{us}{#1s}}

% ---------------------------------------------------------------- interfaces I1-I8
% Nombres cortos EXACTOS (catalogo de la especificacion, seccion 2). Usar siempre estas
% macros en etiquetas y puertos: \itf{3} -> I3 ; \itfn{3} -> I3 observacion y recompensa ;
% \itfnn{3} -> el mismo nombre en dos lineas (I3 observacion / y recompensa), para los
% huecos estrechos (va dentro de un nodo con align, o en una etiqueta con align=center).
% Revision 27.09: I6 = metricas por corrida e I7 = politica publicada (antes al reves), para
% que la arquitectura se lea I3, I4, I6, I7, I8 de arriba abajo.
\newcommand{\nombreinterfaz}[1]{\ifcase#1\relax\or escenario del caso\or
  aproximaciones, enlaces y vecinos\or observación y recompensa\or acción\or
  mensajes entre vecinos\or métricas por corrida\or política publicada\or veredicto\fi}
\newcommand{\nombreinterfazdos}[1]{\ifcase#1\relax\or escenario\\ del caso\or
  aproximaciones,\\ enlaces y vecinos\or observación\\ y recompensa\or acción\or
  mensajes\\ entre vecinos\or métricas\\ por corrida\or política\\ publicada\or veredicto\fi}
\newcommand{\itf}[1]{\textbf{I#1}}
\newcommand{\itfn}[1]{\textbf{I#1}~\nombreinterfaz{#1}}
\newcommand{\itfnn}[1]{\textbf{I#1}~\nombreinterfazdos{#1}}
% texto de un nodo: etiqueta en negrita y detalle debajo
\newcommand{\nodo}[2]{\textbf{#1}\\#2}
% babel spanish pone coma decimal en modo matematico; la tesis escribe punto (f = 2.0)
\AtBeginDocument{\ifdefined\decimalpoint\decimalpoint\fi}

% ---------------------------------------------------------------- estilos
\tikzset{
  % estilo de la figura completa
  % sin cortes de palabra con guion dentro de los nodos (la prueba partia "cru-da")
  diagrama/.style={font=\footnotesize, line join=round, line cap=round,
                   node distance=4mm and 5mm, inner sep=2pt,
                   execute at begin node={\hyphenpenalty=10000\exhyphenpenalty=10000}},
  % ---- agrupaciones (fondo) ------------------------------------------------
  % grupo=<u>: unidad completa; subgrupo=<u>: etapa dentro de una unidad
  grupo/.style={rectangle, rounded corners=5pt, line width=0.8pt, draw=#1b, fill=#1f,
                inner sep=0pt},
  grupo/.default=u,
  subgrupo/.style={rectangle, rounded corners=4pt, line width=0.5pt, draw=#1b, fill=#1s,
                   inner sep=0pt},
  subgrupo/.default=u,
  titulo/.style={font=\small\bfseries, text=#1b, anchor=north west, inner sep=2pt,
                 align=left},
  titulo/.default=u,
  subtitulo/.style={font=\footnotesize\bfseries, text=#1b, anchor=north west,
                    inner sep=2pt, align=left},
  subtitulo/.default=u,
  % ---- nodos --------------------------------------------------------------
  base nodo/.style={draw=#1b, line width=0.6pt, align=center, font=\footnotesize,
                    inner xsep=3pt, inner ysep=2.5pt},
  proceso/.style={base nodo=#1, rectangle, rounded corners=3pt, fill=white,
                  minimum height=7mm},
  proceso/.default=u,
  archivo/.style={base nodo=#1, chamfered rectangle, chamfered rectangle corners=north east,
                  chamfered rectangle xsep=2.2mm, chamfered rectangle ysep=2.2mm,
                  fill=white, minimum height=7mm},
  archivo/.default=u,
  decision/.style={base nodo=#1, diamond, aspect=2.8, fill=white, inner xsep=0.5pt,
                   inner ysep=0.5pt},
  decision/.default=u,
  iniciofin/.style={base nodo=#1, rounded rectangle, fill=#1f, line width=1pt,
                    inner xsep=5pt, minimum height=7mm},
  iniciofin/.default=u,
  % entorno: SUMO + TraCI (U3). Mismo aspecto en los cinco diagramas.
  entorno/.style={rectangle, draw=black!75, double=black!12, double distance=1.2pt,
                  line width=0.5pt, fill=black!12, align=center, font=\footnotesize,
                  inner xsep=4pt, inner ysep=3pt, minimum height=8mm},
  % puerto de interfaz con otra unidad: flecha-etiqueta con el color de ESA unidad.
  % Apunta en el sentido del dato: hacia la derecha por defecto; "hacia izquierda" invierte.
  % Un puerto arriba o abajo del recuadro apunta en vertical ("hacia abajo", "hacia arriba"),
  % con el texto horizontal, igual que su flecha.
  puerto/.style={base nodo=#1, signal, signal to=east, signal pointer angle=120,
                 fill=#1f, line width=0.8pt, align=left, inner xsep=3pt, inner ysep=2pt,
                 minimum height=7mm},
  puerto/.default=u,
  hacia izquierda/.style={signal to=west},
  hacia abajo/.style={signal to=south, signal pointer angle=150, inner ysep=3pt},
  hacia arriba/.style={signal to=north, signal pointer angle=150, inner ysep=3pt},
  % semillas de SUMO de una etapa (Tabla tab:plan2-semillas): "<etapa>: <semillas>"
  semillas/.style={rounded rectangle, draw=black!55, fill=black!5, line width=0.5pt,
                   font=\footnotesize\itshape, align=center, inner xsep=4pt,
                   inner ysep=1.5pt},
  % texto suelto (aclaraciones cortas)
  nota/.style={font=\footnotesize, text=black!85, align=left, inner sep=1pt},
  % modificador: elemento previsto y no usado (se aplica DESPUES del tipo de nodo)
  previsto/.style={draw=black!50, dash pattern=on 0.9pt off 1.5pt, line width=0.8pt,
                   text=black!60, fill=white},
  % marca de aplicabilidad (T, D, P en U2; variantes en U4)
  marca/.style={rectangle, rounded corners=1.5pt, draw=#1b, fill=white, text=#1b,
                font=\footnotesize\bfseries, line width=0.5pt, inner xsep=1.5pt,
                inner ysep=1pt},
  marca/.default=u,
  % barra comun de los aprendices de U2 y U4 (solo en la arquitectura)
  bus/.style={rectangle, rounded corners=1.5pt, draw=black!65, fill=black!10,
              line width=0.6pt, minimum width=4mm, inner sep=0pt},
  % ---- flechas ------------------------------------------------------------
  punta/.style={-{Stealth[length=2.1mm, width=1.6mm]}},
  flujo/.style={punta, draw=black!85, line width=0.6pt, rounded corners=3pt},
  interfaz/.style={punta, draw=black!85, line width=0.6pt, rounded corners=3pt,
                   dash pattern=on 3pt off 1.8pt},
  repeticion/.style={punta, draw=black!85, line width=1.2pt, rounded corners=4pt},
  no usado/.style={punta, draw=black!50, line width=0.8pt, rounded corners=3pt,
                   dash pattern=on 0.9pt off 1.5pt},
  doble/.style={{Stealth[length=2.1mm, width=1.6mm]}-{Stealth[length=2.1mm, width=1.6mm]}},
  % etiqueta de flecha (corta el trazo con fondo blanco; "sin fondo" sobre grupos)
  etiqueta/.style={font=\footnotesize, inner sep=1pt, align=center, fill=white},
  etiqueta sin fondo/.style={font=\footnotesize, inner sep=1pt, align=center},
  etiqueta en grupo/.style={etiqueta, fill=#1f},   % sobre el fondo de un grupo
  etiqueta en grupo/.default=u,
  % ---- leyenda ------------------------------------------------------------
  % Iconos en gris neutro (ext) en las cinco figuras. El archivo toma su tamano de text width y
  % text height (con minimum width/height el biselado vacio no se dibuja), y el texto fija alto y
  % profundidad para que todos queden en la misma linea base.
  ley icono/.style={minimum width=5.5mm, minimum height=3.4mm, inner sep=0pt},
  ley flecha/.style={minimum width=7mm, minimum height=3.4mm, inner sep=0pt,
                     append after command={(\tikzlastnode.west) edge[#1] (\tikzlastnode.east)}},
  ley texto/.style={font=\footnotesize, inner sep=0pt, text height=1.5ex, text depth=.25ex},
  ley/proceso/.style={proceso=ext, ley icono, minimum height=3.4mm},
  ley/archivo/.style={archivo=ext, chamfered rectangle xsep=1.1mm, chamfered rectangle ysep=1.1mm,
                      inner sep=0pt, minimum width=0pt, minimum height=0pt,
                      text width=3.3mm, text height=1.2mm, text depth=0pt},
  ley/decision/.style={decision=ext, minimum width=7mm, minimum height=3.8mm, inner sep=0pt},
  ley/iniciofin/.style={iniciofin=ext, ley icono, minimum height=3.4mm, inner xsep=0pt},
  ley/entorno/.style={entorno, ley icono, minimum height=3.4mm, inner xsep=0pt},
  ley/puerto/.style={puerto=ext, ley icono, minimum height=3.4mm, inner xsep=0pt,
                     minimum width=6.5mm},
  ley/semillas/.style={semillas, ley icono, minimum height=3.4mm, inner xsep=0pt},
  ley/previsto/.style={proceso=ext, previsto, ley icono, minimum height=3.4mm},
  ley/flujo/.style={ley flecha=flujo},
  ley/interfaz/.style={ley flecha=interfaz},
  ley/repeticion/.style={ley flecha=repeticion},
  ley/nousado/.style={ley flecha=no usado},
}
% textos de la leyenda (iguales en los cinco diagramas)
\pgfkeys{/leyenda/.is family, /leyenda,
  proceso/.initial={proceso},
  archivo/.initial={archivo o dato},
  decision/.initial={decisión},
  iniciofin/.initial={inicio o fin},
  entorno/.initial={simulador SUMO},
  puerto/.initial={puerto de interfaz},
  semillas/.initial={semillas},
  previsto/.initial={previsto, no usado},
  flujo/.initial={flujo},
  interfaz/.initial={interfaz (I1--I8)},
  repeticion/.initial={repetición},
  nousado/.initial={camino previsto, no usado},
}
% \leyenda{<punto>}{<claves>}: una fila de leyenda que empieza en <punto> (borde izquierdo,
% altura media). Claves, en este orden: proceso, archivo, decision, iniciofin, entorno,
% puerto, semillas, previsto, flujo, interfaz, repeticion, nousado. Se ponen solo las que
% el diagrama usa; si no caben en 15 cm, dos filas (dos llamadas, 6 mm entre ellas).
% nousado (flecha punteada) se omite si ya esta previsto (el mismo concepto).
\newcommand{\leyenda}[2]{%
  \begin{scope}[start chain=leyenda going right]
    \coordinate[on chain] (leyenda-inicio) at #1;
    \foreach \clave [count=\k] in {#2}{%
      \ifnum\k=1 \def\separacion{0mm}\else \def\separacion{3.5mm}\fi
      \edef\nodoicono{\noexpand\node[node distance=\separacion, on chain, ley/\clave] {};}%
      \nodoicono
      \node[node distance=1.3mm, on chain, ley texto] {\pgfkeysvalueof{/leyenda/\clave}};
    }
  \end{scope}}

% y dibuja a 15 cm exactos de ancho (\useasboundingbox), para que la tesis lo incluya con
% \includegraphics[width=\textwidth] a escala 1. Letra: \footnotesize (8 pt) en nodos y
% etiquetas, \small (9 pt) en los titulos de unidad. Nunca \scriptsize ni \tiny.
% construir.py compila <nombre>.tex y copia el PDF a ../images/diag_<nombre>.pdf.
% Un diagrama se compila aparte: dentro no van \ref ni \cite (se escribe "U2", "Anexo D").
%
% Lenguaje visual
%   color      = unidad a la que pertenece el elemento (U1 azul, U2 naranja, U3 verde,
%                U4 violeta; gris = fuera de las unidades); relleno muy claro y borde
%                oscuro, legible en blanco y negro. Un puerto lleva el color de la unidad
%                del otro lado de la interfaz (gris si son varias).
%   forma      = tipo de nodo: proceso (rectangulo redondeado), archivo (esquina cortada),
%                decision (rombo), inicio/fin (estadio), entorno SUMO + TraCI (gris, doble
%                borde), puerto de interfaz (flecha-etiqueta), semillas (pildora en cursiva).
%   trazo      = tipo de flecha: flujo (continua), interfaz I1-I8 (discontinua: dato de una
%                interfaz del catalogo, lleva \itf o \itfn), repeticion (continua gruesa),
%                no usado (punteada gris; el borde de un nodo previsto es igual).

\usetikzlibrary{positioning, fit, calc, arrows.meta, backgrounds, chains,
                shapes.geometric, shapes.misc, shapes.symbols, babel}

% ---------------------------------------------------------------- colores
% <u>b = borde y titulo; <u>f = fondo del grupo de la unidad; <u>s = fondo de un subgrupo
\definecolor{u1b}{HTML}{2F5597}  \definecolor{u1f}{HTML}{EDF2FA}  \definecolor{u1s}{HTML}{DCE6F5}
\definecolor{u2b}{HTML}{A65A0B}  \definecolor{u2f}{HTML}{FDF3E7}  \definecolor{u2s}{HTML}{F9E3CA}
\definecolor{u3b}{HTML}{3A6E21}  \definecolor{u3f}{HTML}{EEF6EA}  \definecolor{u3s}{HTML}{DDEDD5}
\definecolor{u4b}{HTML}{5E4A9A}  \definecolor{u4f}{HTML}{F2EFF9}  \definecolor{u4s}{HTML}{E3DDF2}
% externo: resultados de la tesis (Cap. V y Anexo A) y todo lo que no es de una unidad
\definecolor{extb}{HTML}{404040} \definecolor{extf}{HTML}{F2F2F2} \definecolor{exts}{HTML}{E6E6E6}
% unidad "actual" (ub, uf, us): la fija \unidad{u2} (o u1, u3, u4, ext) al comienzo de cada
% diagrama; los estilos sin argumento usan esa unidad y los que llevan =<u> usan la indicada
% (p. ej. proceso=u3 para un nodo de U3 dibujado en el diagrama de U2).
\colorlet{ub}{extb} \colorlet{uf}{extf} \colorlet{us}{exts}
\newcommand{\unidad}[1]{\colorlet{ub}{#1b}\colorlet{uf}{#1f}\colorlet{us}{#1s}}

% ---------------------------------------------------------------- interfaces I1-I8
% Nombres cortos EXACTOS (catalogo de la especificacion, seccion 2). Usar siempre estas
% macros en etiquetas y puertos: \itf{3} -> I3 ; \itfn{3} -> I3 observacion y recompensa ;
% \itfnn{3} -> el mismo nombre en dos lineas (I3 observacion / y recompensa), para los
% huecos estrechos (va dentro de un nodo con align, o en una etiqueta con align=center).
% Revision 27.09: I6 = metricas por corrida e I7 = politica publicada (antes al reves), para
% que la arquitectura se lea I3, I4, I6, I7, I8 de arriba abajo.
\newcommand{\nombreinterfaz}[1]{\ifcase#1\relax\or escenario del caso\or
  aproximaciones, enlaces y vecinos\or observación y recompensa\or acción\or
  mensajes entre vecinos\or métricas por corrida\or política publicada\or veredicto\fi}
\newcommand{\nombreinterfazdos}[1]{\ifcase#1\relax\or escenario\\ del caso\or
  aproximaciones,\\ enlaces y vecinos\or observación\\ y recompensa\or acción\or
  mensajes\\ entre vecinos\or métricas\\ por corrida\or política\\ publicada\or veredicto\fi}
\newcommand{\itf}[1]{\textbf{I#1}}
\newcommand{\itfn}[1]{\textbf{I#1}~\nombreinterfaz{#1}}
\newcommand{\itfnn}[1]{\textbf{I#1}~\nombreinterfazdos{#1}}
% texto de un nodo: etiqueta en negrita y detalle debajo
\newcommand{\nodo}[2]{\textbf{#1}\\#2}
% babel spanish pone coma decimal en modo matematico; la tesis escribe punto (f = 2.0)
\AtBeginDocument{\ifdefined\decimalpoint\decimalpoint\fi}

% ---------------------------------------------------------------- estilos
\tikzset{
  % estilo de la figura completa
  % sin cortes de palabra con guion dentro de los nodos (la prueba partia "cru-da")
  diagrama/.style={font=\footnotesize, line join=round, line cap=round,
                   node distance=4mm and 5mm, inner sep=2pt,
                   execute at begin node={\hyphenpenalty=10000\exhyphenpenalty=10000}},
  % ---- agrupaciones (fondo) ------------------------------------------------
  % grupo=<u>: unidad completa; subgrupo=<u>: etapa dentro de una unidad
  grupo/.style={rectangle, rounded corners=5pt, line width=0.8pt, draw=#1b, fill=#1f,
                inner sep=0pt},
  grupo/.default=u,
  subgrupo/.style={rectangle, rounded corners=4pt, line width=0.5pt, draw=#1b, fill=#1s,
                   inner sep=0pt},
  subgrupo/.default=u,
  titulo/.style={font=\small\bfseries, text=#1b, anchor=north west, inner sep=2pt,
                 align=left},
  titulo/.default=u,
  subtitulo/.style={font=\footnotesize\bfseries, text=#1b, anchor=north west,
                    inner sep=2pt, align=left},
  subtitulo/.default=u,
  % ---- nodos --------------------------------------------------------------
  base nodo/.style={draw=#1b, line width=0.6pt, align=center, font=\footnotesize,
                    inner xsep=3pt, inner ysep=2.5pt},
  proceso/.style={base nodo=#1, rectangle, rounded corners=3pt, fill=white,
                  minimum height=7mm},
  proceso/.default=u,
  archivo/.style={base nodo=#1, chamfered rectangle, chamfered rectangle corners=north east,
                  chamfered rectangle xsep=2.2mm, chamfered rectangle ysep=2.2mm,
                  fill=white, minimum height=7mm},
  archivo/.default=u,
  decision/.style={base nodo=#1, diamond, aspect=2.8, fill=white, inner xsep=0.5pt,
                   inner ysep=0.5pt},
  decision/.default=u,
  iniciofin/.style={base nodo=#1, rounded rectangle, fill=#1f, line width=1pt,
                    inner xsep=5pt, minimum height=7mm},
  iniciofin/.default=u,
  % entorno: SUMO + TraCI (U3). Mismo aspecto en los cinco diagramas.
  entorno/.style={rectangle, draw=black!75, double=black!12, double distance=1.2pt,
                  line width=0.5pt, fill=black!12, align=center, font=\footnotesize,
                  inner xsep=4pt, inner ysep=3pt, minimum height=8mm},
  % puerto de interfaz con otra unidad: flecha-etiqueta con el color de ESA unidad.
  % Apunta en el sentido del dato: hacia la derecha por defecto; "hacia izquierda" invierte.
  % Un puerto arriba o abajo del recuadro apunta en vertical ("hacia abajo", "hacia arriba"),
  % con el texto horizontal, igual que su flecha.
  puerto/.style={base nodo=#1, signal, signal to=east, signal pointer angle=120,
                 fill=#1f, line width=0.8pt, align=left, inner xsep=3pt, inner ysep=2pt,
                 minimum height=7mm},
  puerto/.default=u,
  hacia izquierda/.style={signal to=west},
  hacia abajo/.style={signal to=south, signal pointer angle=150, inner ysep=3pt},
  hacia arriba/.style={signal to=north, signal pointer angle=150, inner ysep=3pt},
  % semillas de SUMO de una etapa (Tabla tab:plan2-semillas): "<etapa>: <semillas>"
  semillas/.style={rounded rectangle, draw=black!55, fill=black!5, line width=0.5pt,
                   font=\footnotesize\itshape, align=center, inner xsep=4pt,
                   inner ysep=1.5pt},
  % texto suelto (aclaraciones cortas)
  nota/.style={font=\footnotesize, text=black!85, align=left, inner sep=1pt},
  % modificador: elemento previsto y no usado (se aplica DESPUES del tipo de nodo)
  previsto/.style={draw=black!50, dash pattern=on 0.9pt off 1.5pt, line width=0.8pt,
                   text=black!60, fill=white},
  % marca de aplicabilidad (T, D, P en U2; variantes en U4)
  marca/.style={rectangle, rounded corners=1.5pt, draw=#1b, fill=white, text=#1b,
                font=\footnotesize\bfseries, line width=0.5pt, inner xsep=1.5pt,
                inner ysep=1pt},
  marca/.default=u,
  % barra comun de los aprendices de U2 y U4 (solo en la arquitectura)
  bus/.style={rectangle, rounded corners=1.5pt, draw=black!65, fill=black!10,
              line width=0.6pt, minimum width=4mm, inner sep=0pt},
  % ---- flechas ------------------------------------------------------------
  punta/.style={-{Stealth[length=2.1mm, width=1.6mm]}},
  flujo/.style={punta, draw=black!85, line width=0.6pt, rounded corners=3pt},
  interfaz/.style={punta, draw=black!85, line width=0.6pt, rounded corners=3pt,
                   dash pattern=on 3pt off 1.8pt},
  repeticion/.style={punta, draw=black!85, line width=1.2pt, rounded corners=4pt},
  no usado/.style={punta, draw=black!50, line width=0.8pt, rounded corners=3pt,
                   dash pattern=on 0.9pt off 1.5pt},
  doble/.style={{Stealth[length=2.1mm, width=1.6mm]}-{Stealth[length=2.1mm, width=1.6mm]}},
  % etiqueta de flecha (corta el trazo con fondo blanco; "sin fondo" sobre grupos)
  etiqueta/.style={font=\footnotesize, inner sep=1pt, align=center, fill=white},
  etiqueta sin fondo/.style={font=\footnotesize, inner sep=1pt, align=center},
  etiqueta en grupo/.style={etiqueta, fill=#1f},   % sobre el fondo de un grupo
  etiqueta en grupo/.default=u,
  % ---- leyenda ------------------------------------------------------------
  % Iconos en gris neutro (ext) en las cinco figuras. El archivo toma su tamano de text width y
  % text height (con minimum width/height el biselado vacio no se dibuja), y el texto fija alto y
  % profundidad para que todos queden en la misma linea base.
  ley icono/.style={minimum width=5.5mm, minimum height=3.4mm, inner sep=0pt},
  ley flecha/.style={minimum width=7mm, minimum height=3.4mm, inner sep=0pt,
                     append after command={(\tikzlastnode.west) edge[#1] (\tikzlastnode.east)}},
  ley texto/.style={font=\footnotesize, inner sep=0pt, text height=1.5ex, text depth=.25ex},
  ley/proceso/.style={proceso=ext, ley icono, minimum height=3.4mm},
  ley/archivo/.style={archivo=ext, chamfered rectangle xsep=1.1mm, chamfered rectangle ysep=1.1mm,
                      inner sep=0pt, minimum width=0pt, minimum height=0pt,
                      text width=3.3mm, text height=1.2mm, text depth=0pt},
  ley/decision/.style={decision=ext, minimum width=7mm, minimum height=3.8mm, inner sep=0pt},
  ley/iniciofin/.style={iniciofin=ext, ley icono, minimum height=3.4mm, inner xsep=0pt},
  ley/entorno/.style={entorno, ley icono, minimum height=3.4mm, inner xsep=0pt},
  ley/puerto/.style={puerto=ext, ley icono, minimum height=3.4mm, inner xsep=0pt,
                     minimum width=6.5mm},
  ley/semillas/.style={semillas, ley icono, minimum height=3.4mm, inner xsep=0pt},
  ley/previsto/.style={proceso=ext, previsto, ley icono, minimum height=3.4mm},
  ley/flujo/.style={ley flecha=flujo},
  ley/interfaz/.style={ley flecha=interfaz},
  ley/repeticion/.style={ley flecha=repeticion},
  ley/nousado/.style={ley flecha=no usado},
}
% textos de la leyenda (iguales en los cinco diagramas)
\pgfkeys{/leyenda/.is family, /leyenda,
  proceso/.initial={proceso},
  archivo/.initial={archivo o dato},
  decision/.initial={decisión},
  iniciofin/.initial={inicio o fin},
  entorno/.initial={simulador SUMO},
  puerto/.initial={puerto de interfaz},
  semillas/.initial={semillas},
  previsto/.initial={previsto, no usado},
  flujo/.initial={flujo},
  interfaz/.initial={interfaz (I1--I8)},
  repeticion/.initial={repetición},
  nousado/.initial={camino previsto, no usado},
}
% \leyenda{<punto>}{<claves>}: una fila de leyenda que empieza en <punto> (borde izquierdo,
% altura media). Claves, en este orden: proceso, archivo, decision, iniciofin, entorno,
% puerto, semillas, previsto, flujo, interfaz, repeticion, nousado. Se ponen solo las que
% el diagrama usa; si no caben en 15 cm, dos filas (dos llamadas, 6 mm entre ellas).
% nousado (flecha punteada) se omite si ya esta previsto (el mismo concepto).
\newcommand{\leyenda}[2]{%
  \begin{scope}[start chain=leyenda going right]
    \coordinate[on chain] (leyenda-inicio) at #1;
    \foreach \clave [count=\k] in {#2}{%
      \ifnum\k=1 \def\separacion{0mm}\else \def\separacion{3.5mm}\fi
      \edef\nodoicono{\noexpand\node[node distance=\separacion, on chain, ley/\clave] {};}%
      \nodoicono
      \node[node distance=1.3mm, on chain, ley texto] {\pgfkeysvalueof{/leyenda/\clave}};
    }
  \end{scope}}

% y dibuja a 15 cm exactos de ancho (\useasboundingbox), para que la tesis lo incluya con
% \includegraphics[width=\textwidth] a escala 1. Letra: \footnotesize (8 pt) en nodos y
% etiquetas, \small (9 pt) en los titulos de unidad. Nunca \scriptsize ni \tiny.
% construir.py compila <nombre>.tex y copia el PDF a ../images/diag_<nombre>.pdf.
% Un diagrama se compila aparte: dentro no van \ref ni \cite (se escribe "U2", "Anexo D").
%
% Lenguaje visual
%   color      = unidad a la que pertenece el elemento (U1 azul, U2 naranja, U3 verde,
%                U4 violeta; gris = fuera de las unidades); relleno muy claro y borde
%                oscuro, legible en blanco y negro. Un puerto lleva el color de la unidad
%                del otro lado de la interfaz (gris si son varias).
%   forma      = tipo de nodo: proceso (rectangulo redondeado), archivo (esquina cortada),
%                decision (rombo), inicio/fin (estadio), entorno SUMO + TraCI (gris, doble
%                borde), puerto de interfaz (flecha-etiqueta), semillas (pildora en cursiva).
%   trazo      = tipo de flecha: flujo (continua), interfaz I1-I8 (discontinua: dato de una
%                interfaz del catalogo, lleva \itf o \itfn), repeticion (continua gruesa),
%                no usado (punteada gris; el borde de un nodo previsto es igual).

\usetikzlibrary{positioning, fit, calc, arrows.meta, backgrounds, chains,
                shapes.geometric, shapes.misc, shapes.symbols, babel}

% ---------------------------------------------------------------- colores
% <u>b = borde y titulo; <u>f = fondo del grupo de la unidad; <u>s = fondo de un subgrupo
\definecolor{u1b}{HTML}{2F5597}  \definecolor{u1f}{HTML}{EDF2FA}  \definecolor{u1s}{HTML}{DCE6F5}
\definecolor{u2b}{HTML}{A65A0B}  \definecolor{u2f}{HTML}{FDF3E7}  \definecolor{u2s}{HTML}{F9E3CA}
\definecolor{u3b}{HTML}{3A6E21}  \definecolor{u3f}{HTML}{EEF6EA}  \definecolor{u3s}{HTML}{DDEDD5}
\definecolor{u4b}{HTML}{5E4A9A}  \definecolor{u4f}{HTML}{F2EFF9}  \definecolor{u4s}{HTML}{E3DDF2}
% externo: resultados de la tesis (Cap. V y Anexo A) y todo lo que no es de una unidad
\definecolor{extb}{HTML}{404040} \definecolor{extf}{HTML}{F2F2F2} \definecolor{exts}{HTML}{E6E6E6}
% unidad "actual" (ub, uf, us): la fija \unidad{u2} (o u1, u3, u4, ext) al comienzo de cada
% diagrama; los estilos sin argumento usan esa unidad y los que llevan =<u> usan la indicada
% (p. ej. proceso=u3 para un nodo de U3 dibujado en el diagrama de U2).
\colorlet{ub}{extb} \colorlet{uf}{extf} \colorlet{us}{exts}
\newcommand{\unidad}[1]{\colorlet{ub}{#1b}\colorlet{uf}{#1f}\colorlet{us}{#1s}}

% ---------------------------------------------------------------- interfaces I1-I8
% Nombres cortos EXACTOS (catalogo de la especificacion, seccion 2). Usar siempre estas
% macros en etiquetas y puertos: \itf{3} -> I3 ; \itfn{3} -> I3 observacion y recompensa ;
% \itfnn{3} -> el mismo nombre en dos lineas (I3 observacion / y recompensa), para los
% huecos estrechos (va dentro de un nodo con align, o en una etiqueta con align=center).
% Revision 27.09: I6 = metricas por corrida e I7 = politica publicada (antes al reves), para
% que la arquitectura se lea I3, I4, I6, I7, I8 de arriba abajo.
\newcommand{\nombreinterfaz}[1]{\ifcase#1\relax\or escenario del caso\or
  aproximaciones, enlaces y vecinos\or observación y recompensa\or acción\or
  mensajes entre vecinos\or métricas por corrida\or política publicada\or veredicto\fi}
\newcommand{\nombreinterfazdos}[1]{\ifcase#1\relax\or escenario\\ del caso\or
  aproximaciones,\\ enlaces y vecinos\or observación\\ y recompensa\or acción\or
  mensajes\\ entre vecinos\or métricas\\ por corrida\or política\\ publicada\or veredicto\fi}
\newcommand{\itf}[1]{\textbf{I#1}}
\newcommand{\itfn}[1]{\textbf{I#1}~\nombreinterfaz{#1}}
\newcommand{\itfnn}[1]{\textbf{I#1}~\nombreinterfazdos{#1}}
% texto de un nodo: etiqueta en negrita y detalle debajo
\newcommand{\nodo}[2]{\textbf{#1}\\#2}
% babel spanish pone coma decimal en modo matematico; la tesis escribe punto (f = 2.0)
\AtBeginDocument{\ifdefined\decimalpoint\decimalpoint\fi}

% ---------------------------------------------------------------- estilos
\tikzset{
  % estilo de la figura completa
  % sin cortes de palabra con guion dentro de los nodos (la prueba partia "cru-da")
  diagrama/.style={font=\footnotesize, line join=round, line cap=round,
                   node distance=4mm and 5mm, inner sep=2pt,
                   execute at begin node={\hyphenpenalty=10000\exhyphenpenalty=10000}},
  % ---- agrupaciones (fondo) ------------------------------------------------
  % grupo=<u>: unidad completa; subgrupo=<u>: etapa dentro de una unidad
  grupo/.style={rectangle, rounded corners=5pt, line width=0.8pt, draw=#1b, fill=#1f,
                inner sep=0pt},
  grupo/.default=u,
  subgrupo/.style={rectangle, rounded corners=4pt, line width=0.5pt, draw=#1b, fill=#1s,
                   inner sep=0pt},
  subgrupo/.default=u,
  titulo/.style={font=\small\bfseries, text=#1b, anchor=north west, inner sep=2pt,
                 align=left},
  titulo/.default=u,
  subtitulo/.style={font=\footnotesize\bfseries, text=#1b, anchor=north west,
                    inner sep=2pt, align=left},
  subtitulo/.default=u,
  % ---- nodos --------------------------------------------------------------
  base nodo/.style={draw=#1b, line width=0.6pt, align=center, font=\footnotesize,
                    inner xsep=3pt, inner ysep=2.5pt},
  proceso/.style={base nodo=#1, rectangle, rounded corners=3pt, fill=white,
                  minimum height=7mm},
  proceso/.default=u,
  archivo/.style={base nodo=#1, chamfered rectangle, chamfered rectangle corners=north east,
                  chamfered rectangle xsep=2.2mm, chamfered rectangle ysep=2.2mm,
                  fill=white, minimum height=7mm},
  archivo/.default=u,
  decision/.style={base nodo=#1, diamond, aspect=2.8, fill=white, inner xsep=0.5pt,
                   inner ysep=0.5pt},
  decision/.default=u,
  iniciofin/.style={base nodo=#1, rounded rectangle, fill=#1f, line width=1pt,
                    inner xsep=5pt, minimum height=7mm},
  iniciofin/.default=u,
  % entorno: SUMO + TraCI (U3). Mismo aspecto en los cinco diagramas.
  entorno/.style={rectangle, draw=black!75, double=black!12, double distance=1.2pt,
                  line width=0.5pt, fill=black!12, align=center, font=\footnotesize,
                  inner xsep=4pt, inner ysep=3pt, minimum height=8mm},
  % puerto de interfaz con otra unidad: flecha-etiqueta con el color de ESA unidad.
  % Apunta en el sentido del dato: hacia la derecha por defecto; "hacia izquierda" invierte.
  % Un puerto arriba o abajo del recuadro apunta en vertical ("hacia abajo", "hacia arriba"),
  % con el texto horizontal, igual que su flecha.
  puerto/.style={base nodo=#1, signal, signal to=east, signal pointer angle=120,
                 fill=#1f, line width=0.8pt, align=left, inner xsep=3pt, inner ysep=2pt,
                 minimum height=7mm},
  puerto/.default=u,
  hacia izquierda/.style={signal to=west},
  hacia abajo/.style={signal to=south, signal pointer angle=150, inner ysep=3pt},
  hacia arriba/.style={signal to=north, signal pointer angle=150, inner ysep=3pt},
  % semillas de SUMO de una etapa (Tabla tab:plan2-semillas): "<etapa>: <semillas>"
  semillas/.style={rounded rectangle, draw=black!55, fill=black!5, line width=0.5pt,
                   font=\footnotesize\itshape, align=center, inner xsep=4pt,
                   inner ysep=1.5pt},
  % texto suelto (aclaraciones cortas)
  nota/.style={font=\footnotesize, text=black!85, align=left, inner sep=1pt},
  % modificador: elemento previsto y no usado (se aplica DESPUES del tipo de nodo)
  previsto/.style={draw=black!50, dash pattern=on 0.9pt off 1.5pt, line width=0.8pt,
                   text=black!60, fill=white},
  % marca de aplicabilidad (T, D, P en U2; variantes en U4)
  marca/.style={rectangle, rounded corners=1.5pt, draw=#1b, fill=white, text=#1b,
                font=\footnotesize\bfseries, line width=0.5pt, inner xsep=1.5pt,
                inner ysep=1pt},
  marca/.default=u,
  % barra comun de los aprendices de U2 y U4 (solo en la arquitectura)
  bus/.style={rectangle, rounded corners=1.5pt, draw=black!65, fill=black!10,
              line width=0.6pt, minimum width=4mm, inner sep=0pt},
  % ---- flechas ------------------------------------------------------------
  punta/.style={-{Stealth[length=2.1mm, width=1.6mm]}},
  flujo/.style={punta, draw=black!85, line width=0.6pt, rounded corners=3pt},
  interfaz/.style={punta, draw=black!85, line width=0.6pt, rounded corners=3pt,
                   dash pattern=on 3pt off 1.8pt},
  repeticion/.style={punta, draw=black!85, line width=1.2pt, rounded corners=4pt},
  no usado/.style={punta, draw=black!50, line width=0.8pt, rounded corners=3pt,
                   dash pattern=on 0.9pt off 1.5pt},
  doble/.style={{Stealth[length=2.1mm, width=1.6mm]}-{Stealth[length=2.1mm, width=1.6mm]}},
  % etiqueta de flecha (corta el trazo con fondo blanco; "sin fondo" sobre grupos)
  etiqueta/.style={font=\footnotesize, inner sep=1pt, align=center, fill=white},
  etiqueta sin fondo/.style={font=\footnotesize, inner sep=1pt, align=center},
  etiqueta en grupo/.style={etiqueta, fill=#1f},   % sobre el fondo de un grupo
  etiqueta en grupo/.default=u,
  % ---- leyenda ------------------------------------------------------------
  % Iconos en gris neutro (ext) en las cinco figuras. El archivo toma su tamano de text width y
  % text height (con minimum width/height el biselado vacio no se dibuja), y el texto fija alto y
  % profundidad para que todos queden en la misma linea base.
  ley icono/.style={minimum width=5.5mm, minimum height=3.4mm, inner sep=0pt},
  ley flecha/.style={minimum width=7mm, minimum height=3.4mm, inner sep=0pt,
                     append after command={(\tikzlastnode.west) edge[#1] (\tikzlastnode.east)}},
  ley texto/.style={font=\footnotesize, inner sep=0pt, text height=1.5ex, text depth=.25ex},
  ley/proceso/.style={proceso=ext, ley icono, minimum height=3.4mm},
  ley/archivo/.style={archivo=ext, chamfered rectangle xsep=1.1mm, chamfered rectangle ysep=1.1mm,
                      inner sep=0pt, minimum width=0pt, minimum height=0pt,
                      text width=3.3mm, text height=1.2mm, text depth=0pt},
  ley/decision/.style={decision=ext, minimum width=7mm, minimum height=3.8mm, inner sep=0pt},
  ley/iniciofin/.style={iniciofin=ext, ley icono, minimum height=3.4mm, inner xsep=0pt},
  ley/entorno/.style={entorno, ley icono, minimum height=3.4mm, inner xsep=0pt},
  ley/puerto/.style={puerto=ext, ley icono, minimum height=3.4mm, inner xsep=0pt,
                     minimum width=6.5mm},
  ley/semillas/.style={semillas, ley icono, minimum height=3.4mm, inner xsep=0pt},
  ley/previsto/.style={proceso=ext, previsto, ley icono, minimum height=3.4mm},
  ley/flujo/.style={ley flecha=flujo},
  ley/interfaz/.style={ley flecha=interfaz},
  ley/repeticion/.style={ley flecha=repeticion},
  ley/nousado/.style={ley flecha=no usado},
}
% textos de la leyenda (iguales en los cinco diagramas)
\pgfkeys{/leyenda/.is family, /leyenda,
  proceso/.initial={proceso},
  archivo/.initial={archivo o dato},
  decision/.initial={decisión},
  iniciofin/.initial={inicio o fin},
  entorno/.initial={simulador SUMO},
  puerto/.initial={puerto de interfaz},
  semillas/.initial={semillas},
  previsto/.initial={previsto, no usado},
  flujo/.initial={flujo},
  interfaz/.initial={interfaz (I1--I8)},
  repeticion/.initial={repetición},
  nousado/.initial={camino previsto, no usado},
}
% \leyenda{<punto>}{<claves>}: una fila de leyenda que empieza en <punto> (borde izquierdo,
% altura media). Claves, en este orden: proceso, archivo, decision, iniciofin, entorno,
% puerto, semillas, previsto, flujo, interfaz, repeticion, nousado. Se ponen solo las que
% el diagrama usa; si no caben en 15 cm, dos filas (dos llamadas, 6 mm entre ellas).
% nousado (flecha punteada) se omite si ya esta previsto (el mismo concepto).
\newcommand{\leyenda}[2]{%
  \begin{scope}[start chain=leyenda going right]
    \coordinate[on chain] (leyenda-inicio) at #1;
    \foreach \clave [count=\k] in {#2}{%
      \ifnum\k=1 \def\separacion{0mm}\else \def\separacion{3.5mm}\fi
      \edef\nodoicono{\noexpand\node[node distance=\separacion, on chain, ley/\clave] {};}%
      \nodoicono
      \node[node distance=1.3mm, on chain, ley texto] {\pgfkeysvalueof{/leyenda/\clave}};
    }
  \end{scope}}

% y dibuja a 15 cm exactos de ancho (\useasboundingbox), para que la tesis lo incluya con
% \includegraphics[width=\textwidth] a escala 1. Letra: \footnotesize (8 pt) en nodos y
% etiquetas, \small (9 pt) en los titulos de unidad. Nunca \scriptsize ni \tiny.
% construir.py compila <nombre>.tex y copia el PDF a ../images/diag_<nombre>.pdf.
% Un diagrama se compila aparte: dentro no van \ref ni \cite (se escribe "U2", "Anexo D").
%
% Lenguaje visual
%   color      = unidad a la que pertenece el elemento (U1 azul, U2 naranja, U3 verde,
%                U4 violeta; gris = fuera de las unidades); relleno muy claro y borde
%                oscuro, legible en blanco y negro. Un puerto lleva el color de la unidad
%                del otro lado de la interfaz (gris si son varias).
%   forma      = tipo de nodo: proceso (rectangulo redondeado), archivo (esquina cortada),
%                decision (rombo), inicio/fin (estadio), entorno SUMO + TraCI (gris, doble
%                borde), puerto de interfaz (flecha-etiqueta), semillas (pildora en cursiva).
%   trazo      = tipo de flecha: flujo (continua), interfaz I1-I8 (discontinua: dato de una
%                interfaz del catalogo, lleva \itf o \itfn), repeticion (continua gruesa),
%                no usado (punteada gris; el borde de un nodo previsto es igual).

\usetikzlibrary{positioning, fit, calc, arrows.meta, backgrounds, chains,
                shapes.geometric, shapes.misc, shapes.symbols, babel}

% ---------------------------------------------------------------- colores
% <u>b = borde y titulo; <u>f = fondo del grupo de la unidad; <u>s = fondo de un subgrupo
\definecolor{u1b}{HTML}{2F5597}  \definecolor{u1f}{HTML}{EDF2FA}  \definecolor{u1s}{HTML}{DCE6F5}
\definecolor{u2b}{HTML}{A65A0B}  \definecolor{u2f}{HTML}{FDF3E7}  \definecolor{u2s}{HTML}{F9E3CA}
\definecolor{u3b}{HTML}{3A6E21}  \definecolor{u3f}{HTML}{EEF6EA}  \definecolor{u3s}{HTML}{DDEDD5}
\definecolor{u4b}{HTML}{5E4A9A}  \definecolor{u4f}{HTML}{F2EFF9}  \definecolor{u4s}{HTML}{E3DDF2}
% externo: resultados de la tesis (Cap. V y Anexo A) y todo lo que no es de una unidad
\definecolor{extb}{HTML}{404040} \definecolor{extf}{HTML}{F2F2F2} \definecolor{exts}{HTML}{E6E6E6}
% unidad "actual" (ub, uf, us): la fija \unidad{u2} (o u1, u3, u4, ext) al comienzo de cada
% diagrama; los estilos sin argumento usan esa unidad y los que llevan =<u> usan la indicada
% (p. ej. proceso=u3 para un nodo de U3 dibujado en el diagrama de U2).
\colorlet{ub}{extb} \colorlet{uf}{extf} \colorlet{us}{exts}
\newcommand{\unidad}[1]{\colorlet{ub}{#1b}\colorlet{uf}{#1f}\colorlet{us}{#1s}}

% ---------------------------------------------------------------- interfaces I1-I8
% Nombres cortos EXACTOS (catalogo de la especificacion, seccion 2). Usar siempre estas
% macros en etiquetas y puertos: \itf{3} -> I3 ; \itfn{3} -> I3 observacion y recompensa ;
% \itfnn{3} -> el mismo nombre en dos lineas (I3 observacion / y recompensa), para los
% huecos estrechos (va dentro de un nodo con align, o en una etiqueta con align=center).
% Revision 27.09: I6 = metricas por corrida e I7 = politica publicada (antes al reves), para
% que la arquitectura se lea I3, I4, I6, I7, I8 de arriba abajo.
\newcommand{\nombreinterfaz}[1]{\ifcase#1\relax\or escenario del caso\or
  aproximaciones, enlaces y vecinos\or observación y recompensa\or acción\or
  mensajes entre vecinos\or métricas por corrida\or política publicada\or veredicto\fi}
\newcommand{\nombreinterfazdos}[1]{\ifcase#1\relax\or escenario\\ del caso\or
  aproximaciones,\\ enlaces y vecinos\or observación\\ y recompensa\or acción\or
  mensajes\\ entre vecinos\or métricas\\ por corrida\or política\\ publicada\or veredicto\fi}
\newcommand{\itf}[1]{\textbf{I#1}}
\newcommand{\itfn}[1]{\textbf{I#1}~\nombreinterfaz{#1}}
\newcommand{\itfnn}[1]{\textbf{I#1}~\nombreinterfazdos{#1}}
% texto de un nodo: etiqueta en negrita y detalle debajo
\newcommand{\nodo}[2]{\textbf{#1}\\#2}
% babel spanish pone coma decimal en modo matematico; la tesis escribe punto (f = 2.0)
\AtBeginDocument{\ifdefined\decimalpoint\decimalpoint\fi}

% ---------------------------------------------------------------- estilos
\tikzset{
  % estilo de la figura completa
  % sin cortes de palabra con guion dentro de los nodos (la prueba partia "cru-da")
  diagrama/.style={font=\footnotesize, line join=round, line cap=round,
                   node distance=4mm and 5mm, inner sep=2pt,
                   execute at begin node={\hyphenpenalty=10000\exhyphenpenalty=10000}},
  % ---- agrupaciones (fondo) ------------------------------------------------
  % grupo=<u>: unidad completa; subgrupo=<u>: etapa dentro de una unidad
  grupo/.style={rectangle, rounded corners=5pt, line width=0.8pt, draw=#1b, fill=#1f,
                inner sep=0pt},
  grupo/.default=u,
  subgrupo/.style={rectangle, rounded corners=4pt, line width=0.5pt, draw=#1b, fill=#1s,
                   inner sep=0pt},
  subgrupo/.default=u,
  titulo/.style={font=\small\bfseries, text=#1b, anchor=north west, inner sep=2pt,
                 align=left},
  titulo/.default=u,
  subtitulo/.style={font=\footnotesize\bfseries, text=#1b, anchor=north west,
                    inner sep=2pt, align=left},
  subtitulo/.default=u,
  % ---- nodos --------------------------------------------------------------
  base nodo/.style={draw=#1b, line width=0.6pt, align=center, font=\footnotesize,
                    inner xsep=3pt, inner ysep=2.5pt},
  proceso/.style={base nodo=#1, rectangle, rounded corners=3pt, fill=white,
                  minimum height=7mm},
  proceso/.default=u,
  archivo/.style={base nodo=#1, chamfered rectangle, chamfered rectangle corners=north east,
                  chamfered rectangle xsep=2.2mm, chamfered rectangle ysep=2.2mm,
                  fill=white, minimum height=7mm},
  archivo/.default=u,
  decision/.style={base nodo=#1, diamond, aspect=2.8, fill=white, inner xsep=0.5pt,
                   inner ysep=0.5pt},
  decision/.default=u,
  iniciofin/.style={base nodo=#1, rounded rectangle, fill=#1f, line width=1pt,
                    inner xsep=5pt, minimum height=7mm},
  iniciofin/.default=u,
  % entorno: SUMO + TraCI (U3). Mismo aspecto en los cinco diagramas.
  entorno/.style={rectangle, draw=black!75, double=black!12, double distance=1.2pt,
                  line width=0.5pt, fill=black!12, align=center, font=\footnotesize,
                  inner xsep=4pt, inner ysep=3pt, minimum height=8mm},
  % puerto de interfaz con otra unidad: flecha-etiqueta con el color de ESA unidad.
  % Apunta en el sentido del dato: hacia la derecha por defecto; "hacia izquierda" invierte.
  % Un puerto arriba o abajo del recuadro apunta en vertical ("hacia abajo", "hacia arriba"),
  % con el texto horizontal, igual que su flecha.
  puerto/.style={base nodo=#1, signal, signal to=east, signal pointer angle=120,
                 fill=#1f, line width=0.8pt, align=left, inner xsep=3pt, inner ysep=2pt,
                 minimum height=7mm},
  puerto/.default=u,
  hacia izquierda/.style={signal to=west},
  hacia abajo/.style={signal to=south, signal pointer angle=150, inner ysep=3pt},
  hacia arriba/.style={signal to=north, signal pointer angle=150, inner ysep=3pt},
  % semillas de SUMO de una etapa (Tabla tab:plan2-semillas): "<etapa>: <semillas>"
  semillas/.style={rounded rectangle, draw=black!55, fill=black!5, line width=0.5pt,
                   font=\footnotesize\itshape, align=center, inner xsep=4pt,
                   inner ysep=1.5pt},
  % texto suelto (aclaraciones cortas)
  nota/.style={font=\footnotesize, text=black!85, align=left, inner sep=1pt},
  % modificador: elemento previsto y no usado (se aplica DESPUES del tipo de nodo)
  previsto/.style={draw=black!50, dash pattern=on 0.9pt off 1.5pt, line width=0.8pt,
                   text=black!60, fill=white},
  % marca de aplicabilidad (T, D, P en U2; variantes en U4)
  marca/.style={rectangle, rounded corners=1.5pt, draw=#1b, fill=white, text=#1b,
                font=\footnotesize\bfseries, line width=0.5pt, inner xsep=1.5pt,
                inner ysep=1pt},
  marca/.default=u,
  % barra comun de los aprendices de U2 y U4 (solo en la arquitectura)
  bus/.style={rectangle, rounded corners=1.5pt, draw=black!65, fill=black!10,
              line width=0.6pt, minimum width=4mm, inner sep=0pt},
  % ---- flechas ------------------------------------------------------------
  punta/.style={-{Stealth[length=2.1mm, width=1.6mm]}},
  flujo/.style={punta, draw=black!85, line width=0.6pt, rounded corners=3pt},
  interfaz/.style={punta, draw=black!85, line width=0.6pt, rounded corners=3pt,
                   dash pattern=on 3pt off 1.8pt},
  repeticion/.style={punta, draw=black!85, line width=1.2pt, rounded corners=4pt},
  no usado/.style={punta, draw=black!50, line width=0.8pt, rounded corners=3pt,
                   dash pattern=on 0.9pt off 1.5pt},
  doble/.style={{Stealth[length=2.1mm, width=1.6mm]}-{Stealth[length=2.1mm, width=1.6mm]}},
  % etiqueta de flecha (corta el trazo con fondo blanco; "sin fondo" sobre grupos)
  etiqueta/.style={font=\footnotesize, inner sep=1pt, align=center, fill=white},
  etiqueta sin fondo/.style={font=\footnotesize, inner sep=1pt, align=center},
  etiqueta en grupo/.style={etiqueta, fill=#1f},   % sobre el fondo de un grupo
  etiqueta en grupo/.default=u,
  % ---- leyenda ------------------------------------------------------------
  % Iconos en gris neutro (ext) en las cinco figuras. El archivo toma su tamano de text width y
  % text height (con minimum width/height el biselado vacio no se dibuja), y el texto fija alto y
  % profundidad para que todos queden en la misma linea base.
  ley icono/.style={minimum width=5.5mm, minimum height=3.4mm, inner sep=0pt},
  ley flecha/.style={minimum width=7mm, minimum height=3.4mm, inner sep=0pt,
                     append after command={(\tikzlastnode.west) edge[#1] (\tikzlastnode.east)}},
  ley texto/.style={font=\footnotesize, inner sep=0pt, text height=1.5ex, text depth=.25ex},
  ley/proceso/.style={proceso=ext, ley icono, minimum height=3.4mm},
  ley/archivo/.style={archivo=ext, chamfered rectangle xsep=1.1mm, chamfered rectangle ysep=1.1mm,
                      inner sep=0pt, minimum width=0pt, minimum height=0pt,
                      text width=3.3mm, text height=1.2mm, text depth=0pt},
  ley/decision/.style={decision=ext, minimum width=7mm, minimum height=3.8mm, inner sep=0pt},
  ley/iniciofin/.style={iniciofin=ext, ley icono, minimum height=3.4mm, inner xsep=0pt},
  ley/entorno/.style={entorno, ley icono, minimum height=3.4mm, inner xsep=0pt},
  ley/puerto/.style={puerto=ext, ley icono, minimum height=3.4mm, inner xsep=0pt,
                     minimum width=6.5mm},
  ley/semillas/.style={semillas, ley icono, minimum height=3.4mm, inner xsep=0pt},
  ley/previsto/.style={proceso=ext, previsto, ley icono, minimum height=3.4mm},
  ley/flujo/.style={ley flecha=flujo},
  ley/interfaz/.style={ley flecha=interfaz},
  ley/repeticion/.style={ley flecha=repeticion},
  ley/nousado/.style={ley flecha=no usado},
}
% textos de la leyenda (iguales en los cinco diagramas)
\pgfkeys{/leyenda/.is family, /leyenda,
  proceso/.initial={proceso},
  archivo/.initial={archivo o dato},
  decision/.initial={decisión},
  iniciofin/.initial={inicio o fin},
  entorno/.initial={simulador SUMO},
  puerto/.initial={puerto de interfaz},
  semillas/.initial={semillas},
  previsto/.initial={previsto, no usado},
  flujo/.initial={flujo},
  interfaz/.initial={interfaz (I1--I8)},
  repeticion/.initial={repetición},
  nousado/.initial={camino previsto, no usado},
}
% \leyenda{<punto>}{<claves>}: una fila de leyenda que empieza en <punto> (borde izquierdo,
% altura media). Claves, en este orden: proceso, archivo, decision, iniciofin, entorno,
% puerto, semillas, previsto, flujo, interfaz, repeticion, nousado. Se ponen solo las que
% el diagrama usa; si no caben en 15 cm, dos filas (dos llamadas, 6 mm entre ellas).
% nousado (flecha punteada) se omite si ya esta previsto (el mismo concepto).
\newcommand{\leyenda}[2]{%
  \begin{scope}[start chain=leyenda going right]
    \coordinate[on chain] (leyenda-inicio) at #1;
    \foreach \clave [count=\k] in {#2}{%
      \ifnum\k=1 \def\separacion{0mm}\else \def\separacion{3.5mm}\fi
      \edef\nodoicono{\noexpand\node[node distance=\separacion, on chain, ley/\clave] {};}%
      \nodoicono
      \node[node distance=1.3mm, on chain, ley texto] {\pgfkeysvalueof{/leyenda/\clave}};
    }
  \end{scope}}
